\section{Introduction}
\label{sec:intro}

A central question of equilibrium quantum statistical mechanics is what
survives coarse-graining. Before reduction, equilibrium is generated by a
Hamiltonian and written $\rho\propto e^{-\beta H}$: the state and its generator
are given together. After an environment is traced out, the reduced
equilibrium state remains well defined, but the generator need not remain
transparent. The sharper question is therefore not whether reduced equilibrium
exists, but whether it can still be written as an explicit Hamiltonian rather
than only as an implicit density operator. This is the inside-out version of
equilibrium, where one starts from the reduced state and asks what generator,
if any, survives the reduction.

For an open system with finite coupling, the operational equilibrium state of
the subsystem is the reduced state of the global Gibbs ensemble,
\begin{align}
    \bar{\rho}_Q(\beta) &= \Tr_X\, e^{-\beta H_{\mathrm{tot}}},
    \label{eq:intro_rho_bar} \\
    e^{-\beta H_{\mathrm{MF}}(\beta)}
    &\propto \frac{\Tr_X e^{-\beta H_{\mathrm{tot}}}}{Z_X(\beta)},
    \qquad Z_X(\beta)=\Tr_X e^{-\beta H_X}.
    \label{eq:intro_hmf_def}
\end{align}
The Hamiltonian of mean force (HMF) is the operator that restores a Gibbs-form
description after the bath has been traced out. It is therefore the exact
reduced equilibrium generator whenever system--bath coupling is finite
\cite{campisiFluctuationTheoremArbitrary2009,talknerColloquiumStatisticalMechanics2020,trushechkinOpenQuantumSystem2022,seifertFirstSecondLaw2016,campbellRoadmapQuantumThermodynamics2026,binderThermodynamicsQuantumRegime2018}.

Why care about a mean-force Hamiltonian at all? First, it underlies the
free-energy, work, and entropy relations of strong-coupling thermodynamics
\cite{jarzynskiNonequilibriumWorkTheorem2004,campisiFluctuationTheoremArbitrary2009,seifertFirstSecondLaw2016,millerEntropyProduction2017,lacerdaQuantumThermodynamicsFast2023,BinderVinjanampathyModiGoold2015}.
Second, it provides a consistent reduced equilibrium state when correlations
with the environment cannot be ignored
\cite{pechukasReducedDynamicsNeed1994b,trushechkinOpenQuantumSystem2022,gelinThermodynamicsSubensembleCanonical2009a}.
Third, and more conceptually, it turns a traced Gibbs operator into a
constructive reduced description that can be inspected, truncated, or compared
across operator families.

That constructive viewpoint matters whenever reduced descriptions must remain
explicit, interpretable, or computationally tractable. Recent work on
representation cost emphasizes that coarse-graining is constrained by what a
reduced model class can continue to carry~\cite{mccaulCostComplexity2026}. A
parallel instinct now appears in machine learning through compact
compositional architectures such as Kolmogorov--Arnold networks, which ask
when complicated multivariate dependence can still be carried by a structured
low-complexity decomposition~\cite{liuKANKolmogorovArnoldNetworks2024}. We
return to that comparison only after the physics problem has been made
precise. The present paper remains entirely about equilibrium statistical
mechanics: given an exact reduced Gibbs state, when does it stay inside a
compact operator family?

The literature is broad but structurally clear. Canonical definitions and
thermodynamic identities are now well established
\cite{campisiFluctuationTheoremArbitrary2009,jarzynskiNonequilibriumWorkTheorem2004,seifertFirstSecondLaw2016,talknerColloquiumStatisticalMechanics2020,trushechkinOpenQuantumSystem2022}.
Exact evaluations exist in special cases: commuting interactions close
trivially, quadratic models preserve Gaussianity, and finite-dimensional
examples such as the spin-boson model provide controlled testbeds
\cite{campisiTalknerHanggi2009Solvable,caldeiraQuantumTunnellingDissipative1983a,grabertQuantumBrownianMotion1988,hiltHamiltonianMeanForce2011,leggettDynamicsDissipativeTwostate1987}.
Benchmark asymptotes in weak coupling and in strongly dressed regimes are also
known and remain important checks on any exact construction
\cite{cresserWeakUltrastrongCoupling2021a,hovhannisyanHamiltonianMeanForce2020}.
Recent work has further emphasized operator structure, bath feedback, and
finite-reservoir generalizations
\cite{burkeStructureHamiltonianMean2024,correaPotentialRenormalisationLamb2025,duGeneralizedHamiltonianMeanforce2025a}.

Outside solvable models, however, the HMF is usually accessed numerically or
perturbatively. Reaction-coordinate mappings, polaron transforms, hierarchical
equations, stochastic approaches, and direct imaginary-time solvers can compute
$\bar\rho_Q(\beta)$, but they rarely return a compact closed-form generator
\cite{nazirReactionCoordinateMapping2018,IlesSmithDijkstraLambertNazir2016,AntoSztrikacsNazirSegal2023,ShubrookIlesSmithNazir2025,moixEquilibriumreducedDensityMatrix2012,chenRigorousStochasticMatrix2014,makriExploitingClassicalDecoherence2014,tanimuraReducedHierarchicalEquations2014,songCalculationCorrelatedInitial2015,XuVadimovStockburgerAnkerhold2026,stockburgerSimulatingSpinbosonDynamics2004,mccaulPartitionfreeApproachOpen2017c}.
The obstruction is therefore not the existence of the reduced equilibrium state
but its representability: when does the logarithm of a traced exponential stay
inside a restricted operator family such as few-body terms, local operators, or
a chosen algebra?

The influence-functional formalism is a natural language for this question. For
Gaussian baths with linear coupling it provides an exact route to integrating
out environmental degrees of freedom
\cite{feynmanTheoryGeneralQuantum1963a,caldeiraQuantumTunnellingDissipative1983a,grabertQuantumBrownianMotion1988}.
In equilibrium it becomes an imaginary-time influence functional, which admits
a Hubbard--Stratonovich rewriting as a quenched stochastic average
\cite{hubbardCalculationPartitionFunctions1959a,stratonovich1957QDistro,stockburgerExactNumberRepresentation2002}.
This formulation makes the reduction step explicit: the partial trace becomes
an average over imaginary-time back-action histories, while all remaining
questions are pushed onto the system operator algebra.

The present paper takes the quenched-density construction of
Ref.~\cite{mccaulMeanForceHamiltoniansInfluence2026} as its exact starting
point and uses it to answer a representability question. For Gaussian baths
with linear coupling, bath statistics collapse to scalar kernel moments while
operator growth is generated entirely by the adjoint chain
$\{\mathrm{ad}_{H_Q}^n(f)\}_{n\geq 0}$. In the Gaussian sector the exact HMF
already comes with a compositional architecture: an adjoint orbit, a bilinear
aggregation weighted by bath moments, and a nonlinear BCH recombination.
Conditions (C1)--(C3) are precisely the closure conditions for those three
layers to remain inside a chosen operator ansatz. We then instantiate this
criterion in the spin-boson qubit, where the reduced equilibrium compresses to
a closed Bloch-plane generator controlled by two response channels and a single
influence magnitude. To our knowledge this yields the first closed-form HMF for
a genuinely noncommuting open quantum system. The same variables organize the
numerical crossover, revealing a representability bottleneck once the dressed
reduced state outgrows the chosen finite basis. The KAN comparison then
becomes a statement about the compositional structure already present in the
exact HMF construction itself.
